% When an organization adopts a large language model, it has to decide which model to use and how to use it. The natural move is to consult a benchmark leaderboard, but almost without exception, benchmarks often focus on a single dimension, automation: how well does the model complete a task on its own? But this is rarely how models are used in practice at work. A model far more often augments work, by assisting a person, structuring the problem, surfacing a constraint, or sharpening a draft, while the human ultimately produces the final deliverable.  This gap—between the automation that benchmarks reward and augmentation that organizations actually deploy is the what we examine in this paper.

When an organization adopts a large language model (LLM), it faces two decisions: which model to use, and how to use it. For the first, the natural move is to consult a benchmark leaderboard—but almost without exception, leaderboards measure a single dimension, automation: how well a model completes a task on its own. This is rarely how models are used at work. Far more often a model augments work, assisting a person by structuring the problem, surfacing a constraint, or sharpening a draft, while the worker produces the final deliverable. The gap between the automation that benchmarks reward and the augmentation that organizations actually deploy is what we examine in this paper. As we will show, the best player is not always the best coach.

Computer scientists have built large-scale benchmarks to evaluate and rank how well LLMs complete tasks: MMLU \citep{hendrycks2021} tests knowledge across 57 subjects, HELM \citep{liang2023} evaluates 30 models on 42 scenarios, and a growing ecosystem of professional benchmarks now tests agents on everything from industrial workflows \citep{sopbench2025} to occupational tasks drawn from 65 specialized domains \citep{occubench2026}. Economists and management scholars have run field experiments to measure AI's productivity effects: ChatGPT cut professional writing time by 40\% and raised quality by 18\% \citep{noy2023}, GitHub Copilot cut programming time by more than half \citep{peng2023}, and a large-scale deployment in customer support raised productivity by 15\% \citep{brynjolfsson2025}. Both traditions have produced genuine insight. But neither fully addresses how workers and organizations actually want to use AI. In practice, a model rarely produces a finished deliverable on its own. Instead it assists another individual. For example, a nationwide survey of 1,500 U.S.\ workers across 104 occupations found that ``equal partnership'' with AI - not replacement - was the dominant preference in nearly half of all occupational categories \citep{shao2025}. If benchmarks only measure what a model can do alone, and field experiments only measure whether a specific tool helped in a specific setting, we are missing a big part of the picture.

Therefore, there is a gap that captures the distinction between \emph{automation} and \emph{augmentation}. A model that performs well when asked to complete a task end-to-end may not be the same model that best helps a human worker complete that same task. The best players are not always the best coaches. For example, \citet{agrawal2023} argue that one worker's automation creates augmentation opportunities for another, \citet{fugener2026} formalize the distinction and show that the optimal mode depends on the type of complementarity between human and AI capability, and \citet{marguerit2025} finds at the macro level that augmentation AI raises wages for high-skilled workers while automation AI displaces low-skilled ones - but this has not been built into the way we evaluate models. Field experiments that test augmentation are narrow by design: \citet{dellacqua2023} test one model with BCG consultants, \citet{chen2024} test ghostwriting versus sounding-board modes in advertising, \citet{doshi2024} study AI-assisted creative writing, \citet{ju2026} study human-AI versus human-human advertising teams. Each is informative but none tell you which of today's models is best suited to augmenting workers across a range of professional tasks - because none was designed to. They test a single model, in a single domain, against a single baseline.



This leaves organizations and individuals in a difficult position. When selecting an AI tool for human-in-the-loop work, ``which model is best?'' is the wrong question; ``best for what role, on what kind of task?'' is the right one. \citet{vaccaro2024}, in a meta-analysis of 106 experiments, find that human-AI combinations frequently underperform the best of either alone - a result that makes no sense if you assume the best automator is also the best collaborator, but makes perfect sense if it is not. \citet{haupt2025} make the methodological implication explicit: evaluation-by-imitation rewards systems that replace people rather than systems that augment them, and the field needs evaluations where human and AI solve tasks together. GDPval \citep{gdpval2025} and related benchmarks \citep{profbench2025, jobbench2026} have moved toward economically grounded tasks and pairwise evaluation, and \citet{collins2024} and \citet{chang2025} have shown that AI-alone accuracy does not predict interactive performance. 

Industry has begun to move in this direction as well. The closest prior effort is Upwork's Human+Agent Productivity Index \citep{upwork_hapi2025}, which evaluates frontier agents on roughly 300
real, paid client projects and finds that pairing agents with expert
freelancers raises completion rates by up to 70\% relative to agents
alone. HAPI is important evidence that agent-only scoring understates
economic value, but it differs from our study in three ways. First, it
studies a single human--AI pairing rather than augmentation
\emph{across} models: it contrasts agent-alone with human-in-the-loop,
not model A versus model B as assistants to a fixed worker. Second, in HAPI a human expert assists the AI—reviewing and correcting the agent's output—so the human occupies the coaching role and the model is the party being helped. Our framework reverses this. The model under test is the assistant, providing scaffolding to a fixed worker, so we evaluate models in the coaching role rather than as the worker being coached. This is precisely the capability our analysis foregrounds—whether a model is a good coach, not only a good player. Third, HAPI's coverage is bounded by what can be sourced and staffed on a freelance marketplace--each task requires recruiting a qualified expert--while our fixed-worker design scales to any professional task we can specify, without freelancer recruitment. HAPI thus motivates our setting; we answer the model-selection question it leaves open.

All in all, no existing framework jointly evaluates models as automators and as augmenters across multiple professional domains---the cross-model, cross-mode comparison that would actually inform tool selection does not exist.

To our knowledge, we are the first to evaluate LLMs jointly as automators and augmenters across multiple professional domains. Our pipeline evaluates a set of frontier LLMs across two axes simultaneously: usage mode (automation versus augmentation) and task domain. In automation mode, each model completes a professional task end-to-end. In augmentation mode, each model provides coaching to a fixed worker model, which then completes the task. Holding the worker constant isolates the value of the assistant's guidance from the executor's own capability, producing a direct measure of coaching quality rather than team outcome. This worker is uses GPT 3.5 Turbo to simulate a worker, which allows us to approximate centaur evaluations—pairings of a high-capability assistant with a fixed, lower-capability worker model. This methodological approach allows us to measure differences in augmentation and automation to any model-task pairing without a seperate field experiment.


Both outputs are evaluated by a panel of LLM judges using task-specific rubrics across seven heterogeneous professional domains, with model families excluded from judging their own outputs (e.g. Claude Opus 4.8 cannot judge other Claude model outputs). The result is not only a single leaderboard score but also a cross-regime rank matrix: for each model, we can ask whether it ranks higher as a direct solver or as a planning coach, and whether that ordering holds across tasks. We find model rankings diverge meaningfully across modes and across task domains—the best automator is not always the best augmenter, and no single model dominates across augmentation tasks. Organizations can use this pipeline directly: by running a task typical of their work, they can identify the model best suited to their specific use case, whether that means delegating entirely or keeping humans in the loop. Concretely, the strongest direct solver is not the strongest assistant, and the best augmenting model changes from task to task, so a single leaderboard cannot tell an organization which model to deploy in which role.


